\documentclass{article}
\usepackage[margin=1.15in]{geometry}
\usepackage{xcolor}
\usepackage{parskip}
\usepackage{fancyhdr}
\usepackage{array}
\usepackage{booktabs}
\usepackage{enumitem}
\usepackage{amsmath}
\usepackage{tabularx}
\usepackage{listings}

\pagestyle{fancy}
\fancyhf{}
\cfoot{\thepage}
\rhead{Lecture 18}
\lhead{MA 214: Applied Statistics (Spring 2026)}
\renewcommand{\headrulewidth}{.4mm}

\newcommand{\blankline}[1]{\underline{\hspace{#1}}}

\lstset{
  basicstyle=\ttfamily\small,
  columns=fullflexible,
  frame=single,
  framerule=0.4pt,
  breaklines=true,
  showstringspaces=false
}

\begin{document}

\noindent\textbf{Name:}\blankline{2.25in}\hfill\textbf{BUID:}\blankline{1.55in}

\medskip

\noindent\textbf{Name:}\blankline{2.25in}\hfill\textbf{BUID:}\blankline{1.55in}\newline

\begin{center}
 \Large In-Class Worksheet: Conjugate Models
\end{center}

\subsection*{Scenario}
A \textbf{community health clinic} is reviewing its adult patient population. The director wants to combine what published studies already say with what the clinic actually sees, and has three questions --- one about a \textbf{proportion}, one about a \textbf{mean}, one about a \textbf{rate}:

\begin{enumerate}[label={\bfseries \arabic*.},topsep=2pt,itemsep=2pt]
	\item What fraction of adult patients have \emph{uncontrolled hypertension}?
	\item What is the \emph{mean systolic blood pressure} (SBP, mmHg) among adult patients?
	\item How many \emph{walk-in patients} arrive per hour on the afternoon shift?
\end{enumerate}

\medskip
\noindent\textbf{Goals:}
\textbf{(1)} apply each of the three conjugate update rules,
\textbf{(2)} identify the prior sample size $n_0$ in each,
\textbf{(3)} verify that the posterior mean really is the weighted average the lecture promised, and
\textbf{(4)} say which side --- prior or data --- each answer leans toward, and why.

\medskip
\noindent\textbf{Setup.} You will work in PAIRS. One person is the \textbf{analyst} (writes the code), the other is the \textbf{guide} (reads the questions and fills out the worksheet). Open RStudio and start a new script --- you will write the code yourself.

\subsection*{Part 1: A Proportion --- the Beta-Binomial}
\textbf{\color{red} Do not use Gen AI when completing this part!}

Published studies of similar clinics suggest an uncontrolled-hypertension rate around \textbf{10\%}. The director is not very confident, so she wants a prior worth about \textbf{10 observations}.

\begin{enumerate}[label={\bfseries (\Alph*)},itemsep=.6em]
	\item \textbf{Choose the prior.} Using $E(p) = \frac{\alpha}{\alpha+\beta}$ and the reading of $\alpha + \beta$ as the prior sample size, find $\alpha$ and $\beta$ giving prior mean $0.10$ and $n_0 = 10$. Show your reasoning. \\[4em]

	$\alpha = $ \blankline{2cm} \hspace{1cm} $\beta = $ \blankline{2cm} \hspace{1cm} $n_0 = $ \blankline{2cm}

	\item \textbf{Update.} The clinic reviews $n = 50$ adult charts; $x = 8$ show uncontrolled hypertension. Apply the Beta-Binomial rule.

	\vspace{0.3em}
	$\alpha_{\text{post}} = \alpha + x = $ \blankline{2cm} \hspace{1cm} $\beta_{\text{post}} = \beta + n - x = $ \blankline{2cm}

	\vspace{0.5em}
	Posterior mean $= \dfrac{\alpha_{\text{post}}}{\alpha_{\text{post}} + \beta_{\text{post}}} = $ \blankline{3cm} \hspace{1cm} MLE $= x/n = $ \blankline{2cm}

	\item \textbf{Check the pattern.} Compute the two weights and verify the posterior mean by hand:
	\[
	\underbrace{\tfrac{n_0}{n_0+n}}_{=\;\blankline{1.5cm}} \times \underbrace{(\text{prior mean})}_{=\;\blankline{1.5cm}}
	\;+\;
	\underbrace{\tfrac{n}{n_0+n}}_{=\;\blankline{1.5cm}} \times \underbrace{(\text{MLE})}_{=\;\blankline{1.5cm}}
	\;=\; \blankline{2.5cm}
	\]

	Does it match your answer to (B)? \\[2em]

	\item Which side does the posterior lean toward, and why? \\[3em]

\end{enumerate}

\newpage

\subsection*{Part 2: A Mean --- the Normal-Normal}

Published studies put the mean SBP for adults in primary care at about \textbf{120 mmHg}; the director's confidence corresponds to a prior SD of $\sigma_0 = 10$ mmHg. The clinic records SBP for $n = 20$ patients and finds $\bar{x} = 126$ mmHg. Take the population SD as known: $\sigma = 15$ mmHg.

\begin{enumerate}[label={\bfseries (\Alph*)},itemsep=.6em]
	\item \textbf{The prior.} \qquad $\mu \sim N\left(\blankline{1.5cm},\ \blankline{1.5cm}^2\right)$

	\item \textbf{Prior sample size.} A different family, so a different recipe for $n_0$:
	\[
	n_0 = \frac{\sigma^2}{\sigma_0^2} = \frac{\blankline{1.2cm}^2}{\blankline{1.2cm}^2} = \blankline{2cm}
	\]

	In plain words: this prior is worth about \blankline{2cm} patient measurements.

	\item \textbf{The posterior.}
	\[
	n_0 + n = \blankline{2.5cm}, \qquad \sigma_n = \frac{\sigma}{\sqrt{n_0+n}} = \blankline{2.5cm}
	\]
	\[
	\mu_n = \frac{n_0}{n_0+n}\,\mu_0 + \frac{n}{n_0+n}\,\bar{x} = \hspace{6cm}
	\]

	\item \textbf{Interval.} The posterior is Normal, so an interval reads straight off it. Use \texttt{qnorm()} (or $\mu_n \pm 1.96\,\sigma_n$):

	95\% credible interval: [\blankline{2cm}, \blankline{2cm}]

	\vspace{0.3em}
	Interpret it in one sentence, in context. \\[3em]

	\item \textbf{Compare with Part 1.} There, $n_0 = 10$ against $n = 50$. Here, $n_0 = 2.25$ against $n = 20$. In which of the two analyzes does the prior have more say? Justify with the weights, not intuition. \\[4em]

\end{enumerate}

\newpage

\subsection*{Part 3: A Rate --- the Gamma-Poisson}

Afternoon walk-ins have historically averaged about \textbf{4 per hour}. The director encodes that as a $\text{Gamma}(8, 2)$ prior. Over six afternoon hours the clinic counts:

\begin{center}
\renewcommand{\arraystretch}{1.2}
\begin{tabular}{lcccccc}
\toprule
Hour & 1 & 2 & 3 & 4 & 5 & 6 \\
\midrule
Walk-ins & 6 & 4 & 7 & 5 & 8 & 3 \\
\bottomrule
\end{tabular}
\end{center}

\begin{enumerate}[label={\bfseries (\Alph*)},itemsep=.6em]
	\item \textbf{Summarize the data.} The Poisson likelihood only needs two numbers from the sample --- which two?

	\vspace{0.3em}
	$n = $ \blankline{1.5cm} \hspace{1cm} $\sum x_i = $ \blankline{1.5cm} \hspace{1cm} $\bar{x} = $ \blankline{2cm}

	\item \textbf{Check the prior.} Confirm that $\text{Gamma}(8,2)$ really does encode ``about 4 per hour'': prior mean $= \alpha/\beta = $ \blankline{2cm}

	\vspace{0.3em}
	And its prior sample size is $n_0 = \beta = $ \blankline{1.5cm}

	\item \textbf{Update.}
	\[
	\lambda \mid \mathbf{x} \sim \text{Gamma}\big(\alpha + \textstyle\sum x_i,\ \beta + n\big) = \text{Gamma}\big(\blankline{1.5cm},\ \blankline{1.5cm}\big)
	\]

	Posterior mean $=$ \blankline{2.5cm} \hspace{1cm} Posterior SD $= \sqrt{\alpha_{\text{post}}/\beta_{\text{post}}^2} = $ \blankline{2.5cm}

	\item \textbf{The pattern, a third time.} Weights: \blankline{1.5cm} on the prior, \blankline{1.5cm} on the data.

	\vspace{0.3em}
	Verify: \blankline{1.2cm} $\times$ \blankline{1.2cm} $+$ \blankline{1.2cm} $\times$ \blankline{1.2cm} $=$ \blankline{2cm}

	\item \textbf{Answer the director's question.} She will add an afternoon receptionist if the rate has genuinely risen above 4 per hour. Compute $P(\lambda > 4 \mid \mathbf{x})$ with \texttt{pgamma()}.

	$P(\lambda > 4 \mid \mathbf{x}) = $ \blankline{2.5cm}

	\vspace{0.3em}
	What do you tell her? Note that you answered her question \emph{directly}, with no null hypothesis and no p-value --- say in one sentence why that is possible here but was not in Module 2. \\[4em]

\end{enumerate}

\newpage

\subsection*{Part 4: One Pattern, Three Families}
\textbf{\color{red} Do not use Gen AI when completing this part!}

\begin{enumerate}[label={\bfseries (\Alph*)},itemsep=.6em]
	\item \textbf{Fill in the table} from your work above. This is the whole lecture on one line each.

	\vspace{0.3em}
	\begin{center}
	\renewcommand{\arraystretch}{1.5}
	{\small
	\begin{tabular}{lccccc}
	\toprule
	& \textbf{Parameter} & \textbf{$n_0$ recipe} & \textbf{$n_0$} & \textbf{$n$} & \textbf{Weight on data} \\
	\midrule
	Part 1 (Beta-Bin.)  & $p$       & $\alpha+\beta$          & \blankline{1.2cm} & \blankline{1.2cm} & \blankline{1.8cm} \\
	Part 2 (Normal-Nor.)& $\mu$     & $\sigma^2/\sigma_0^2$   & \blankline{1.2cm} & \blankline{1.2cm} & \blankline{1.8cm} \\
	Part 3 (Gamma-Pois.)& $\lambda$ & $\beta$                 & \blankline{1.2cm} & \blankline{1.2cm} & \blankline{1.8cm} \\
	\bottomrule
	\end{tabular}}
	\end{center}

	\item Rank the three analyzes from ``prior matters most'' to ``prior matters least.'' Which single number decides the ranking? \\[3em]

	\item \textbf{Doubling the data.} Suppose the clinic doubles each sample --- 100 charts, 40 SBP readings, 12 hours of walk-ins --- and every summary statistic ($x/n$, $\bar x$, $\bar x$) comes out identical. Without recomputing anything, say what happens to each posterior mean and to each posterior SD, and why. \\[5em]

	\item \textbf{Where it breaks.} The director now asks: ``what fraction of patients have uncontrolled hypertension \emph{and} arrive as walk-ins \emph{and} are over 65?'' --- three unknowns at once, tangled together. Explain in two sentences why none of today's three update rules helps, and what that tells you about how special conjugacy is. \\[5em]

\end{enumerate}

\end{document}
